\section{Model Cards and Datasheets}

\begin{frame}{Documentation Is Part of the Model}

    A model that ships without documentation forces every downstream user to rediscover its limits
    --- usually in production.

    \vspace{0.8em}

    \begin{columns}[T]
        \begin{column}{0.48\textwidth}
            \begin{tcolorbox}[colback=raiBlue!5, colframe=raiBlue!60, title=Model card]
                \small
                Travels \textbf{with the trained model}. Answers: what is it for, who is it for,
                where does it fail, and how well does it work \emph{per group}.

                \vspace{0.4em}
                Mitchell et al., FAT* 2019.
            \end{tcolorbox}
        \end{column}
        \begin{column}{0.48\textwidth}
            \begin{tcolorbox}[colback=raiGreen!5, colframe=raiGreen!60, title=Datasheet]
                \small
                Travels \textbf{with the dataset}. Answers: why was it collected, from whom, with
                what consent, how was it cleaned, and what should it not be used for.

                \vspace{0.4em}
                Gebru et al., CACM 2021.
            \end{tcolorbox}
        \end{column}
    \end{columns}

    \vspace{1em}

    \begin{block}{Both are borrowed ideas}
        \centering
        \small A model card is a \textbf{datasheet for a component}, as in electronics; a dataset
        datasheet is the same idea one step upstream. Neither is a research artefact --- they are
        engineering hygiene, and increasingly they are what a customer or a partner asks to see
        before they will build on your model.
    \end{block}
\end{frame}

\begin{frame}{What Goes in a Model Card}

    \centering
    \scriptsize
    \begin{tabular}{p{3.2cm} p{8.6cm}}
        \toprule
        \textbf{Section} & \textbf{Content} \\
        \midrule
        Model details      & Owner, date, version, model type, training algorithm, licence, citation, contact. \\
        Intended use       & Primary intended uses and users; \textbf{out-of-scope} uses stated explicitly. \\
        Factors            & Groups, instrumentation, and environments the performance is expected to vary over. \\
        Metrics            & Which measures, which decision threshold, and how uncertainty was estimated. \\
        Evaluation data    & Which datasets, why they were chosen, and how they were preprocessed. \\
        Training data      & Distribution over the same factors, as far as it can be disclosed. \\
        Quantitative analyses & Results \textbf{unitary} (per factor) and \textbf{intersectional} (per combination). \\
        Ethical considerations & Sensitive data, risks to people, mitigations applied, known problematic uses. \\
        Caveats \& recommendations & Untested groups, open concerns, what testing should happen next. \\
        \bottomrule
    \end{tabular}

    \vspace{0.9em}
    \raggedright
    {\scriptsize M.~Mitchell, S.~Wu, A.~Zaldivar, P.~Barnes, L.~Vasserman, B.~Hutchinson, E.~Spitzer,
    I.~D.~Raji and T.~Gebru, ``Model Cards for Model Reporting,'' FAT* 2019, Section 4,
    \href{https://arxiv.org/abs/1810.03993v2}{arXiv:1810.03993v2}.}
\end{frame}
